\begin{Proposition}{次数的可加性}\label{prop:degree-additivity}
  若 $R$ 是整环, 且 $f,g\in R[x_{1},\ldots,x_{n}]$ 非零, 则
  \[
    \deg(fg)=\deg(f)+\deg(g).
  \]
\end{Proposition}

\begin{Proof}
  设 $f,g\in R[x_{1},\ldots,x_{n}]$ 是次数分别为 $d$ 和 $e$ 的非零多项式. 写成
  \[
    f=f_{d}+f_{d-1}+\cdots+f_{0}
    \quad\text{和}\quad
    g=g_{e}+g_{e-1}+\cdots+g_{0},
  \]
  其中 $f_{i}$ 由 $f$ 中次数为 $i$ 的所有项组成, $g_{j}$ 类似. 由假设, $f_{d}\ne0$ 且 $g_{e}\ne0$.
  稍加思考便可看出, 单项式的次数具有可加性. 因此, $fg$ 中可能以非零系数出现的最高次单项式的次数为 $d+e$, 且来自乘积 $f_{d}g_{e}$.
  由于 $R[x_{1},\ldots,x_{n}]$ 是整环, 可知 $f_{d}g_{e}\ne0$, 从而得出 $\deg(fg)=d+e=\deg(f)+\deg(g)$.
\end{Proof}

\BookExerciseHeading
\begin{BookExercises}
  \item\label{ex:polynomial-ring-laws} 回顾环 (交换且含幺) 的定义, 并证明: 对任意环 $R$, 多项式的加法与乘法使 $R[x_{1},\ldots,x_{n}]$ 具有环 (交换且含幺) 的结构.
  \item\label{ex:regroup-three-variable-polynomial} 对多项式
    \[
      f=xyz^{2}+xyz+z^{3}+x^{2}z^{2}+yz^{2}+z+x+1\in R[x,y,z]
    \]
    的各项分组, 将其分别看作 $R[x,y][z]$, $R[x][y,z]$ 和 $R[y,z][x]$ 的元素.
  \item\label{ex:recursive-ring-isomorphism} 证明 Proposition~\ref{prop:polynomial-ring-recursion} 的证明中定义的函数
    \[
      \phi:R[x_{1},\ldots,x_{n}]\to R[x_{1},\ldots,x_{n-1}][x_{n}]
    \]
    是环同构.
  \item\label{ex:constant-polynomials-subring} 若只要 $\alpha\ne(0,\ldots,0)$ 就有 $a_{\alpha}=0$, 则称多项式 $f=\sum_{\alpha}a_{\alpha}x^{\alpha}\in R[x_{1},\ldots,x_{n}]$ 为常数多项式. 证明常数多项式的集合构成 $R[x_{1},\ldots,x_{n}]$ 中与 $R$ 同构的一个子环.
  \item\label{ex:evaluation-homomorphism} 设 $a=(a_{1},\ldots,a_{n})\in R^{n}$. 证明存在唯一的环同态
    \[
      \phi_{a}:R[x_{1},\ldots,x_{n}]\to R
    \]
    使得对所有 $r\in R$ 有 $\phi_{a}(r)=r$, 对所有 $i$ 有 $\phi_{a}(x_{i})=a_{i}$. (这个同态称为在 $a$ 处的求值, $\phi_{a}(f)$ 常写作 $f(a_{1},\ldots,a_{n})$.)
  \item\label{ex:failure-without-integral-domain} 说明若不假设 $R$ 是整环, Propositions~\ref{prop:polynomials-over-domain} 和 \ref{prop:degree-additivity} 不成立.
  \item\label{ex:degree-zero-polynomial} 说明无论怎样选取 $\deg(0)\in\mathbb{N}$, Proposition~\ref{prop:degree-additivity} 都不能推广到这一情形. (有些数学家采用约定 $\deg(0)=-\infty$.)
  \item\label{ex:degree-of-sum} 证明 $\deg(f+g)\leq\max\{\deg(f),\deg(g)\}$.
\end{BookExercises}
